% 06-conclusion — v2.0
We presented a reproducible degradation benchmark for single-channel and low-cost EEG
and applied it to two open datasets. The results are conditional and, we think, useful.
Alpha reactivity survived the acquisition chain almost intact once the reference was
taken at the same electrode (99.99\% retained; a ratio-of-ratios outcome, so gain-like
steps cancel by construction), whereas the electrode-site decision cost
about 42\% of eyes-closed alpha power and sensor noise reduced the effect monotonically,
to about 0.59 of the reference at 0\,dB. Affective valence was decodable from a full high-end
montage ($0.69$) but was not detected from a single electrode ($0.52$); arousal survived
weakly ($0.56$ above chance but far below $0.68$), and four-class accuracy fell to $0.34$ (above chance but far below the montage). A
low-cost 14-channel headset performed at a similar level to the high-end montage without
passing a formal equivalence test. Decomposition showed that the acquisition chain was
harmless, that valence was not distinguishable from chance at any single site and the
four-class result stayed far below the montage, and that beta-band
features carried part of the full-montage advantage --- a result that cautions against
reading the montage number as purely cortical.

The deliverable is a decision matrix, a single-channel reporting checklist, and an open,
seeded analysis pipeline that turns public data into the reported numbers. The next
version requires the full subject sample, a real single-channel device, a validated
artifact model, registration of the analysis plan, and a larger equivalence study. The
value of the exercise is that it replaces a plausible story ("single channels lose
information") with a measured one: what a single electrode can and cannot show depends
on the analysis, the site and the band, and that boundary can be mapped.
